% ═══════════════════════════════════════════════════════════════
% BEAMER-TEMPLATE: Spanisch-Unterricht
% Folien für Present/Apply-Stundenstruktur
% ═══════════════════════════════════════════════════════════════
%
% VERWENDUNG:
% 1. Kopiere diese Datei in dein Projektverzeichnis
% 2. Passe METADATEN an (Titel, Untertitel, Klasse)
% 3. Erstelle Folien nach Present/Apply-Schema:
%    - PRESENT: Inhalt vermitteln
%    - APPLY: → AB Aufgabe X (Praxis-Verweis)
% 4. Kompiliere 2x mit pdflatex
%
% FARBKONVENTION:
% - \textcolor{richtig}{...} → Grün für Haken/Korrektes
% - \textcolor{falsch}{...}  → Rot für Kreuz/Fehler
% - \textcolor{akzent}{...}  → Karminrot für Betonung
% - \stw{...}                → Rot für Stammwechsel
%
% ═══════════════════════════════════════════════════════════════

\documentclass[aspectratio=169, 11pt]{beamer}

\usepackage[utf8]{inputenc}
\usepackage[T1]{fontenc}
\usepackage[ngerman]{babel}

% --- SCHRIFTEN ---
\usepackage{helvet}
\renewcommand{\familydefault}{\sfdefault}
\usepackage{palatino}
\newcommand{\seriftext}[1]{{\fontfamily{ppl}\selectfont #1}}

\usepackage{tikz}
\usepackage{array}
\usepackage{booktabs}
\usepackage{pifont}

% ═══════════════════════════════════════════════════════════════
% FARBEN
% ═══════════════════════════════════════════════════════════════

\definecolor{akzent}{HTML}{C41E3A}      % Karminrot (Spanisch-Fachfarbe)
\definecolor{stamm}{HTML}{C62828}       % Stammwechsel-Highlight
\definecolor{hellgrau}{HTML}{F5F5F5}
\definecolor{dunkelgrau}{HTML}{333333}
\definecolor{mittelgrau}{HTML}{666666}
\definecolor{richtig}{HTML}{2A7A4B}     % Grün für Haken
\definecolor{falsch}{HTML}{C62828}      % Rot für Kreuz

% ═══════════════════════════════════════════════════════════════
% BEAMER-THEME
% ═══════════════════════════════════════════════════════════════

\usetheme{default}
\usecolortheme{default}

% Keine Navigation
\setbeamertemplate{navigation symbols}{}

% Farben
\setbeamercolor{frametitle}{fg=white, bg=akzent}
\setbeamercolor{title}{fg=white}
\setbeamercolor{subtitle}{fg=white}
\setbeamercolor{author}{fg=white}
\setbeamercolor{date}{fg=white}
\setbeamercolor{structure}{fg=akzent}
\setbeamercolor{itemize item}{fg=akzent}
\setbeamercolor{itemize subitem}{fg=mittelgrau}
\setbeamercolor{block title}{fg=white, bg=akzent}
\setbeamercolor{block body}{bg=hellgrau}

% Aufzählungszeichen
\setbeamertemplate{itemize item}{\textbullet}
\setbeamertemplate{itemize subitem}{--}

% Titel-Folie (vollflächig Karminrot)
\setbeamertemplate{title page}{
    \begin{tikzpicture}[remember picture, overlay]
        \fill[akzent] (current page.north west) rectangle (current page.south east);
    \end{tikzpicture}
    \vspace{2cm}
    \begin{center}
        {\Huge\bfseries\textcolor{white}{\inserttitle}}\\[0.8cm]
        {\Large\textcolor{white}{\insertsubtitle}}\\[1.5cm]
        {\large\textcolor{white}{\insertauthor}}\\[0.3cm]
        {\textcolor{white}{\insertdate}}
    \end{center}
}

% Frame-Titel (Karminrot-Balken)
\setbeamertemplate{frametitle}{
    \nointerlineskip
    \begin{beamercolorbox}[wd=\paperwidth, ht=1.2cm, dp=0.3cm]{frametitle}
        \hspace{0.5cm}\Large\bfseries\insertframetitle
    \end{beamercolorbox}
}

% Fußzeile mit Seitenzahl
\setbeamertemplate{footline}{
    \hfill\textcolor{mittelgrau}{\footnotesize\insertframenumber}\hspace{0.5cm}\vspace{0.3cm}
}

% ═══════════════════════════════════════════════════════════════
% BEFEHLE
% ═══════════════════════════════════════════════════════════════

% Spanisch fett (Serif)
\newcommand{\es}[1]{\seriftext{\textbf{#1}}}

% Deutsche Übersetzung (kursiv, grau)
\newcommand{\de}[1]{\textit{\textcolor{mittelgrau}{#1}}}

% Stammwechsel (rot, fett)
\newcommand{\stw}[1]{\textcolor{stamm}{\textbf{#1}}}

% AB-Verweis für Apply-Folien
\newcommand{\abmark}[1]{{\Large$\rightarrow$ \textbf{Arbeitsblatt Aufgabe #1}}}

% Richtig/Falsch-Symbole
\newcommand{\richmark}{\textcolor{richtig}{\ding{51}}}
\newcommand{\falschmark}{\textcolor{falsch}{\ding{55}}}

% ═══════════════════════════════════════════════════════════════
% METADATEN (ANPASSEN!)
% ═══════════════════════════════════════════════════════════════

\title{Thema der Stunde}
\subtitle{Untertitel / Kurzbeschreibung}
\author{Spanisch · Klasse X}
\date{}

% ═══════════════════════════════════════════════════════════════
% DOKUMENT
% ═══════════════════════════════════════════════════════════════

\begin{document}

% ---------------------------------------------------------------
% TITELFOLIE
% ---------------------------------------------------------------
\begin{frame}[plain]
    \titlepage
\end{frame}

% ---------------------------------------------------------------
% LERNZIELE
% ---------------------------------------------------------------
\begin{frame}{Lernziele}
    \vspace{0.5cm}

    {\large Heute lernst du:}

    \vspace{0.8cm}

    \begin{itemize}
        \setlength{\itemsep}{0.6cm}
        \item \textbf{Lernziel 1}\\
              {\small Kurze Erklärung}
        \item \textbf{Lernziel 2}\\
              {\small Kurze Erklärung}
        \item \textbf{Lernziel 3}\\
              {\small Kurze Erklärung}
    \end{itemize}
\end{frame}

% ═══════════════════════════════════════════════════════════════
% BLOCK 1: PRESENT → APPLY
% ═══════════════════════════════════════════════════════════════

% --- PRESENT: Inhalt vermitteln ---
\begin{frame}{Thema 1}
    \vspace{0.3cm}

    % Inhalt hier

\end{frame}

% --- APPLY: Praxis-Verweis ---
\begin{frame}{Jetzt du!}
    \vspace{1cm}

    \begin{center}
        \abmark{1}

        \vspace{1cm}

        {\large Bearbeite die Aufgabe allein / mit Partner.}

        \vspace{0.5cm}

        {\small Zeit: X Minuten}
    \end{center}
\end{frame}

% ═══════════════════════════════════════════════════════════════
% BLOCK 2: PRESENT → APPLY
% ═══════════════════════════════════════════════════════════════

% --- PRESENT ---
\begin{frame}{Thema 2}
    \vspace{0.3cm}

    % Inhalt hier

\end{frame}

% --- APPLY ---
\begin{frame}{Jetzt du!}
    \vspace{1cm}

    \begin{center}
        \abmark{2}

        \vspace{1cm}

        {\large Bearbeite die Aufgabe.}
    \end{center}
\end{frame}

% ═══════════════════════════════════════════════════════════════
% ABSCHLUSS
% ═══════════════════════════════════════════════════════════════

% --- ZUSAMMENFASSUNG ---
\begin{frame}{Zusammenfassung}
    \vspace{0.3cm}

    \begin{columns}[T]
        \begin{column}{0.45\textwidth}
            \textbf{1. Punkt}

            \begin{itemize}
                \item Detail
                \item Detail
            \end{itemize}
        \end{column}

        \begin{column}{0.45\textwidth}
            \textbf{2. Punkt}

            \begin{itemize}
                \item Detail
                \item Detail
            \end{itemize}
        \end{column}
    \end{columns}
\end{frame}

% --- TYPISCHE FEHLER (optional) ---
\begin{frame}{Typische Fehler vermeiden}
    \vspace{0.5cm}

    \begin{center}
        \begin{tabular}{@{}cl@{\hspace{2cm}}cl@{}}
            \textcolor{falsch}{\Large$\times$} & Falscher Satz & \textcolor{richtig}{\Large$\checkmark$} & Richtiger Satz \\
        \end{tabular}
    \end{center}
\end{frame}

\end{document}
